\section{R5: General recursion from incomplete examples}\label{sec:r5}

\question{R5.Q1}{Does a capability interpreter agree with the realized recursive term definitionally?}
\status{Not in general. Propositional simulation is the robust contract, especially for well-founded recursion.}

Lean's documentation explicitly distinguishes structural recursion from
well-founded recursion: well-founded definitions typically compute through
propositional equation theorems rather than convenient definitional reduction,
and kernel evaluation can be slow because it unfolds termination evidence
\cite{leanrec}. Therefore the search interpreter should not promise that
its execution trace is literally the kernel's conversion trace.

Define a typed body language with recursive calls admitted only through a
capability
\[
 \mathsf{rec}:(y:X)\to R(y,x)\to Y(y),
\]
where a separate certificate establishes that $R$ is well-founded. The body
may call the capability at a provably smaller argument; the proof of decrease
is part of the obligation, not an unchecked annotation. Realization constructs
a structural or well-founded definition, and a theorem connects the body
semantics to its equation lemmas.

A practical interpreter is fuelled and returns either a value or unknown:
\[
 \mathsf{eval}_m(b,x)=\mathsf{some}(y)
 \quad\Longrightarrow\quad
 \mathsf{realize}(b)(x)=y.
 \tag{7}\label{eq:simulation}
\]
Prove this theorem by induction on the successful evaluation derivation,
using the realization's equation theorem at each recursive step. Fuel
exhaustion supplies no negative information. A natural-valued decrease measure
may bound call depth, but not the entire number of reductions or the size of
a branching recursion tree; define fuel precisely.

There are three evaluation profiles. In the narrowest profile, kernel
reduction validates every observation directly. In a proof-producing profile,
a fast interpreter supplies certificates used to rewrite the finite contract
to decidable ground obligations. In a native profile, compiled evaluation may
supply hints or rely on explicitly approved native-computation assumptions.
Only the first two preserve a kernel-only computational trust boundary without
additional assumptions. An axiom audit must reflect whichever profile is
actually used; one should not describe native verification as if it were
ordinary kernel reduction \cite{leanvalidation}.

The proposal's compilation defect reinforces this separation. A logical
recursor term can pass the kernel yet fail the code generator's expected
representation. Keep the certified term, then construct a compilable equation
form and prove its correspondence. A proof of logical correspondence and an
explicit compiler trust policy are separate from successfully attaching an
implementation annotation. Realization and certification must be charged to
the reported end-to-end budget, not deferred beyond the advertised time to
solution.

\question{R5.Q2}{What changes when angelic execution is combined with top-down search?}
\status{The combination is coherent, but top-down search needs persistent relational assumptions and must revisit them when shared holes are filled.}

Bottom-up synthesis has complete candidate fragments whose behavior can be
summarized compositionally. A top-down skeleton can contain the same hole in
several contexts, with type assignments still changing. Its angelic summary
is therefore a relation on hole closures, recursive calls, and expected
outputs, not just a set of values for one subtree.

Maintain an assumption store per branch. Give a recursive call a stable key
containing the function being synthesized, typed argument values or symbolic
arguments, and relevant parameters. Equal call keys share constraints. When
unevaluation derives requirements for an unobserved call, record them in the
store; when the call becomes executable, discharge them or refine the branch.
When a shared hole changes, invalidate dependent assumptions as well as
forward evaluations.

The safe initial version uses these assumptions to prioritize holes and
candidate bodies but never prunes solely because one guessed execution fails.
A stronger version may prune when a certified relational over-approximation
has no satisfying assignment. This preserves the basic existential meaning
of angelic execution and the conservative-refutation condition of
Section~\ref{sec:semantic}. Burst supplies the main precedent, while Smyth
shows the importance of interdependent goals and live evaluation
\cite{burst,smyth}.

Recent work also broadens the comparison beyond examples. Cataclyst, described
in a July 2026 preprint, combines sketches, enumerative search, mixed-quantifier
first-order properties, and counterexample-derived syntactic constraints.
Its setting and verification infrastructure differ from Lean's, but it is a
relevant source of refinements for recursive synthesis with relational
contracts \cite{cataclyst}. A counterexample-guided restriction imported from
such a system still needs a Lean-side explanation of why it cannot eliminate
a valid candidate in the represented grammar.

\question{R5.Q3}{Is there a small canonical basis of recursion schemes with measured coverage?}
\status{There are effective published template bases, notably Para, but no unique small basis whose reported coverage transfers automatically to Leant2.}

A paramorphism exposes both a child's original input and its recursively
computed result. For a polynomial functor $F$, its algebra has shape
\[
 \phi:F(\mu F\times K)\to K,
\]
and the recursive computation replaces each child by the pair of that child
and its recursive result before applying $\phi$. An ordinary fold is the
special case that ignores the original child. This gives one structured
interface for many apparently different recursive programs.

Para demonstrates synthesis using such templates and recursion-free hole
fillers, including representations useful for multi-argument and higher-order
computations \cite{para}. Leant2 should adapt the template language rather
than mechanically copy a surface syntax. The appropriate starting basis is
structural recursion/paramorphism, bounded product and function-valued state,
bounded nesting, and an explicit well-founded adapter. Tupled state handles
Fibonacci; a function-valued state handles accumulator passing; one structural
argument often suffices for zip; merge can decrease the sum of two list
lengths; Euclidean gcd needs a modulus-decrease theorem, not merely an
unconditional call to an arithmetic tactic.

A basis can be small because it hides large search problems in its algebras
or termination proofs. Thus ``one universal recursor'' is not a useful answer
to the coverage question. Measure both scheme coverage with ideal fillers and
actual synthesis coverage under the same component inventory and resource
budget. Missing comparisons, append, or arithmetic lemmas can otherwise be
mistaken for missing recursion schemes.

Mutual and nested inductives should be handled through recursor metadata and
Lean's equation elaborator, not brittle rules based on generated declaration
names. Initially admit a clearly documented subset and return unsupported for
other forms. Helpers should be synthesized as parameterized subproblems with
explicit call interfaces, avoiding hidden access to a parent's changing local
context. A bound of one helper is an experiment, not an expressiveness theorem.
The roadmap should retain both carrier-given and carrier-invented benchmark
tracks for every recursive family.
